% =====================================================================
%  第 9 章：临界水平 ĝ-模的结构（章首 + §9.1–9.3）
%  原文 PDF p.261–p.283（印刷 p.251–273，PDF − 10）
%  片段止于 PDF p.283 页末一句中部，下接 ch09-2.tex
% =====================================================================

\chapter{临界水平 $\gghat$-模的结构}

在前几章里，我们描述了临界水平顶点代数 $V_{\kappa_c}(\fg)$ 的中心 $\mathfrak{z}(\gghat)$。如 3.3.2 节所证，$\mathfrak{z}(\gghat)$ 同构于 $\gghat_{\kappa_c}$-模 $V_{\kappa_c}(\fg)$ 的、与 $\gghat_{\kappa_c}$ 的作用对易的自同态代数。我们已把这个代数等同于圆盘 $D$ 上 ${}^L G$-opers 的空间 $\Op_{{}^L G}(D)$ 上的函数代数。\transnote{卡点：临界水平为何特殊。背景：$\gghat_{\kappa_c}$ 的分量由 Sugawara 构造给出 Virasoro 算子，其表达式含因子 $(k+h^\vee)^{-1}$；在 $k=-h^\vee$（即 $\kappa_c=-h^\vee\kappa_0$，$\kappa_0$ 为归一化不变内积）处该因子发散，通常的顶点代数结构退化，而 $V_{\kappa_c}(\fg)$ 的中心——即与 $\gghat_{\kappa_c}$ 作用对易的自同态全体——反而变大。结论：本章处理的正是这一临界水平，「中心变大」是它最直接的后果。[来源：原书 p.251 与 §3.3.2；Kac, Infinite-Dimensional Lie Algebras, 3rd ed., CUP 1990（Sugawara 构造一章）]}

本章对临界水平 Verma 模与 Weyl 模的自同态代数得到类似结果。我们证明，二者都是完备泛包络代数 $\widetilde{U}_{\kappa_c}(\gghat)$ 的中心 $Z(\gghat)$ 的商；而该中心同构于去心圆盘 $D^\times$ 上 ${}^L G$-opers 的空间 $\Op_{{}^L G}(D^\times)$ 上的函数代数。在 Verma 模的情形，自同态代数等同于带正则奇点且留数固定的 opers 空间上的函数代数；在 Weyl 模的情形，则是带正则奇点、留数固定且单值性平凡的 opers 空间上的函数代数。于是，opers 的几何反映在临界水平仿射代数的表示论之中。理解表示论与几何之间的这一联系，对发展下一章将要讨论的局部 Langlands 对应十分重要。

本章先引入 $\Op_{{}^L G}(D^\times)$ 中相关的子空间：带正则奇点的 opers（9.1 节）与幂零 opers（9.2 节），并说明它们之间的相互关系。然后在 9.3 节考虑 8.2.1 节定义的 Miura 变换在带正则奇点的 opers 与幂零 opers 上的限制。我们把幂零 opers 上 Miura 变换的纤维与 Springer 纤维联系起来。这将使我们得到以 Springer 纤维为参数族的 Wakimoto 模族。接着在 9.4 节讨论“生活”在带正则奇点 opers 空间之上的若干 $\gghat_{\kappa_c}$-表示范畴。最后在 9.5、9.6 节描述临界水平 Verma 模与 Weyl 模的自同态代数。

9.1 节的结果属于 A. Beilinson 与 V. Drinfeld [1997]；本章其余各节的大部分结果由 D. Gaitsgory 与我在 [2006c; 2007c] 中得到。

\begin{keypoints}
    \item \textbf{核心思想：}临界水平 $\gghat$-模的自同态代数由 opers 的几何控制；本章把 Verma 模、Weyl 模分别与“带正则奇点 + 固定留数”“…… + 单值性平凡”的 opers 空间对应起来。
    \item \textbf{数学工具：}opers 的典范代表（Kostant 截面 $\fg^{\mathrm{can}}$）、幂零 opers、Deligne 扩张、Miura 变换与其纤维、Springer 纤维。
    \item \textbf{关键结论：}$\Op^{RS}_G(D)\subset\Op_G(D^\times)$（命题 9.1.1）；幂零 opers 与带正则奇点 opers 的同构（命题 9.2.1）；Miura 变换在幂零 oper 上的纤维是留数的 Springer 纤维（式 (9.3-9)–(9.3-11)、定理 9.3.6）。
    \item \textbf{记号提醒：}$\check\lambda,\check\mu,\check\rho,\check\omega_i$ 是 $\fh$ 中带“check”的量（对偶侧），$\varpi$ 是投影 $\fh\to\fh/W$；$h$ 既表示 Cartan 子代数也表示 Coxeter 数，按上下文区分。
\end{keypoints}

\section{带正则奇点的 opers}

\begin{keypoints}
    \item \textbf{核心思想：}把圆盘 $D$ 上“在原点有正则奇点”的 oper 定义为带形式参数 $1/t$ 的联络的规范等价类，并用留数记录其首项。
    \item \textbf{数学工具：}规范变换、Kostant 的典范截面 $\fg^{\mathrm{can}}\cong\fh/W$、坐标无关的射影联络与高阶微分表述。
    \item \textbf{关键结论：}命题 9.1.1：$\Op^{RS}_G(D)$ 单射入 $\Op_G(D^\times)$，像是典范代表形如 (9.1-5)、留数由 (9.1-6) 给出的 opers；推论 9.1.2 给出坐标无关的同构。
    \item \textbf{记号提醒：}$\Op^{RS}_G(D)_\varpi$ 表示留数类为 $\varpi$ 的子空间；$E=\{d_1,\dots,d_\ell\}$ 是 $\fg$ 的指数（计重数）。
\end{keypoints}

本节按 Beilinson 与 Drinfeld [1997] 引入圆盘上在原点有正则奇点的 opers 空间。

\subsection{定义}

回忆 $D$（相应地 $D^\times$）上 $G$-opers 的空间 $\Op_G(D)$（相应地 $\Op_G(D^\times)$）是形如 (4.2-1) 的算子空间商去 $B[[t]]$（相应地 $B((t))$）的作用，其中 $\psi_i(t)$ 与 $\mathbf{v}(t)$ 取值于 $\C[[t]]$（相应地 $\C((t))$）。

圆盘 $D$ 上带正则奇点的 $G$-oper，按定义（见 [Beilinson and Drinfeld 1997] 第 3.8.8 节）是形如
\begin{equation}
\nabla = \partial_t + t^{-1}\left(\sum_{i=1}^{\ell}\psi_i(t)f_i + \mathbf{v}(t)\right),
\tag{9.1-1}
\end{equation}
的算子的 $B[[t]]$-共轭类，其中 $\psi_i(t)\in\C[[t]]$，$\psi_i(0)\ne 0$，$\mathbf{v}(t)\in\fb[[t]]$。

等价地，它是形如
\begin{equation}
\nabla = \partial_t + \frac{1}{t}\left(p_{-1} + \mathbf{v}(t)\right),\qquad \mathbf{v}(t)\in\fb[[t]]
\tag{9.1-2}
\end{equation}
的算子的 $N[[t]]$-等价类。

把 $D$ 上带正则奇点的 opers 空间记作 $\Op^{RS}_G(D)$。

更一般地，按 [Beilinson and Drinfeld 1997] 第 3.8.8 节，我们考虑奇点阶数不超过 $k$（$k>0$）的 opers 空间 $\Op^{\mathrm{ord}\,k}_G(D)$，即形如
\[
\nabla = \partial_t + \frac{1}{t^k}\left(p_{-1} + \mathbf{v}(t)\right),\qquad \mathbf{v}(t)\in\fb[[t]]
\]
的算子的 $N[[t]]$-等价类构成的空间。

\subsection{留数}

按 [Beilinson and Drinfeld 1997]，我们给带正则奇点的 oper 配上其\textbf{留数}。对形如 (9.1-2) 的算子，留数按定义等于 $p_{-1}+\mathbf{v}(0)$。显然，在 $N[[t]]$ 的元素 $x(t)$ 作规范变换时，留数被 $x(0)\in N$ 共轭。因此它在
\[
\fg/G = \Spec(\operatorname{Fun}\fg)^G = \Spec(\operatorname{Fun}\fh)^W = \fh/W
\]
中的投影良定义。于是得到态射
\[
\operatorname{res}: \Op^{RS}_G(D)\to \fh/W.
\]
给定 $\varpi\in\fh/W$，把带正则奇点且留数为 $\varpi$ 的 opers 空间记作 $\Op^{RS}_G(D)_\varpi$。

\subsection{典范代表}

设给定一个带正则奇点的 oper，它由形如 (9.1-1) 的联络 $\nabla$ 的 $B[[t]]$-共轭类代表。回忆集合 $E=\{d_1,\dots,d_\ell\}$ 是 $\fg$ 的指数（计重数），而 4.2.4 节引入了 $\fg$ 的带基 $p_1,\dots,p_\ell$ 的子空间 $V^{\mathrm{can}} = \bigoplus_{i\in E}V^{\mathrm{can}}_i$。

与引理 4.2.2 的证明相同，我们利用 $B[[t]]$ 的规范作用把 $\nabla$ 化为
\begin{equation}
\partial_t + t^{-1}\left(p_{-1} + \sum_{j\in E}v_j(t)p_j\right),\qquad v_j(t)\in\C[[t]].
\tag{9.1-3}
\end{equation}
这个算子的留数等于
\begin{equation}
p_{-1} + \sum_{j\in E}v_j(0)p_j.
\tag{9.1-4}
\end{equation}
把 $\fg$ 中由一切形如 $p_{-1}+\sum_{j\in E}y_jp_j$ 的元素构成的仿射子空间记作 $\fg^{\mathrm{can}}$。回忆 [Kostant 1963]：李代数 $\fg$ 中任一正则元素的伴随轨道含唯一的属于 $\fg^{\mathrm{can}}$ 的元素，且相应的态射 $\fg^{\mathrm{can}}\to\fh/W$ 是同构。于是可知：当我们把一个带正则奇点的 oper 化成典范形 (9.1-3) 时，其留数就实现为 $\fg^{\mathrm{can}}$ 中的一个元素。\transnote{卡点：$p_{-1}$ 的下标 $-1$ 与留数良定义的关系。背景：$p_{-1}=\sum_{i=1}^{\ell}x_i^-$ 是最高根 $-\theta$ 对应的根向量之和，与 $p_1$、$2\check\rho$ 构成主 $\mathfrak{sl}_2$-三元组；Kostant 定理断言正则半单元素的伴随轨道与仿射子空间 $\fg^{\mathrm{can}}$ 恰交于一点。结论：正因如此，(9.1-3) 中常数项 $p_{-1}+\sum_{j\in E}v_j(0)p_j$ 在 $\fh/W$ 中的类不依赖 Borel 规范的选择，留数才良定义。[来源：原书 p.253；Kostant 1963, Amer. J. Math. 85, 327--404]}

% [Tutor Note: 这里 p_{−1} 的下标是 −1（不是 1），它对应 $\mathfrak{sl}_2$ 生成元 f 的推广；Kostant 定理保证“正则伴随轨道 ∩ g^can 恰为一点”，这正是留数被良定义为 h/W 中元素的根据。抄写时极易把 p_{−1} 写成 p_1，务必注意。]

考虑自然态射 $\Op^{RS}_G(D)\to\Op_G(D^\times)$，它把形如 (9.1-1) 的算子的 $B[[t]]$-等价类映到其 $B((t))$-等价类。

\begin{proposition}[{[Drinfeld and Sokolov 1985], Prop. 3.8.9}]
映射 $\Op^{RS}_G(D)\to\Op_G(D^\times)$ 是单射。它的像由 $D^\times$ 上那些典范代表形如
\begin{equation}
\partial_t + p_{-1} + \sum_{j\in E}t^{-j-1}c_j(t)p_j,\qquad c_j(t)\in\C[[t]]
\tag{9.1-5}
\end{equation}
的 $G$-opers 组成。此外，这个 oper 的留数等于
\begin{equation}
p_{-1} + \left(c_1(0)+\frac{1}{4}\right)p_1 + \sum_{j\in E,\,j>1}c_j(0)p_j.
\tag{9.1-6}
\end{equation}
\end{proposition}

\begin{proof}
先把带正则奇点的 oper 化为 (9.1-3)。接着作 $\check\rho(t)^{-1}$ 的规范变换，得到
\[
\partial_t + p_{-1} + \check\rho t^{-1} + \sum_{j\in E}t^{-j-1}v_j(t)p_j,\qquad v_j(t)\in\C[[t]].
\]
最后作 $\exp(-p_1/2t)$ 的规范变换，得到算子
\begin{equation}
\partial_t + p_{-1} + t^{-2}\left(v_1(t)-\frac{1}{4}\right)p_1 + \sum_{j\in E,\,j>1}t^{-j-1}v_j(t)p_j,\qquad v_j(t)\in\C[[t]].
\tag{9.1-7}
\end{equation}
于是得到带正则奇点的 opers 空间与 $D^\times$ 上形如 (9.1-5) 的 opers 空间之间的同构；特别地，映射 $\Op^{RS}_G(D)\to\Op_G(D^\times)$ 是单射。此外，比较公式 (9.1-7) 与 (9.1-5) 可知 $c_1(t) = v_1(t)-\frac14$，且当 $j>1$ 时 $c_j(t)=v_j(t)$。由 (9.1-4)，由公式 (9.1-5) 给出的 oper 的留数等于 (9.1-6)。\transnote{卡点：$\frac14$ 与 $c_1(t)=v_1(t)-\frac14$ 从何而来。背景：证明末步的规范变换取 $g=\exp(-p_1/2t)$，它一方面经 $g\partial_tg^{-1}$ 改变 $t^{-2}$ 项的系数，另一方面经 $\operatorname{Ad}_g$ 把 $t^{-1}p_1$ 的一部分并入，两项合并后恰留下 $-\frac14$。结论：同一常数在推论 9.1.2 中表现为 $\eta_{1,-2}+\frac14$，两处必须同时成立——只改一处，$p_1$ 的系数就对不上。[来源：原书 p.253--254]}
\end{proof}

% [Tutor Note: $\check\rho(t)^{-1}$ 是“对偶 Weyl 向量”对应的余特征标，$\exp(-p_1/2t)$ 一步产生 (9.1-6) 中关键的 1/4 常数（射影联络的 Schwarz 导数项）。这个 1/4 是本节最容易抄错的地方之一，它与 9.3 节 Miura 变换中的 $\check\rho$ 平移互相呼应。]

作为推论，对 $\fh/W\simeq\fg^{\mathrm{can}}$ 中任一点 $\varpi$，带正则奇点且留数为 $\varpi$ 的 opers 构成 $\Op^{RS}_G(D)$ 的仿射子空间。

利用同构 (4.2-9)，可以把命题 9.1.1 表述成坐标无关的形式。把 $D^\times$ 上在原点极点阶数不超过 2 的射影联络空间记作 $\mathcal{P}roj_{\ge-2}$，把 $D^\times$ 上在原点极点阶数不超过 $d_j+1$ 的 $(d_j+1)$-微分空间记作 $\omega^{\otimes(d_j+1)}_{\ge-(d_j+1)}$。

\begin{corollary}
$\Op^{RS}_G(D)$ 与空间
\[
\mathcal{P}roj_{\ge-2}\times\bigoplus_{j>1}\omega^{\otimes(d_j+1)}_{\ge-(d_j+1)}
\]
同构。在此同构下，$\Op^{RS}_G(D)_\varpi$ 对应由满足下述条件的 $\ell$-元组 $(\eta_1,\dots,\eta_\ell)$ 构成的仿射子空间：把 $\eta_i$ 按 $t$ 的幂展开的 $t^{-d_i-1}$-系数记作 $\eta_{i,-d_i-1}$（它与 $t$ 的选取无关）。则
\[
p_{-1} + \left(\eta_{1,-2}+\frac{1}{4}\right)p_1 + \sum_{j>1}\eta_{j,-j-1}p_j
\]
在 $\fg/G = \fh/W$ 中的共轭类等于 $\varpi$。
\end{corollary}

设给定一个正则 oper
\[
\partial_t + p_{-1} + \mathbf{v}(t),\qquad \mathbf{v}(t)\in\fb[[t]].
\]
用 $B((t))$ 中元素 $\check\rho(t)$ 作规范变换，把它化为
\[
\partial_t + \frac{1}{t}\left(p_{-1}-\check\rho + t\cdot\check\rho(t)(\mathbf{v}(t))\check\rho(t)^{-1}\right).
\]
若 $\mathbf{v}(t)$ 是正则的，则 $\check\rho(t)\mathbf{v}(t)\check\rho(t)^{-1}$ 也正则。因此这个 oper 具有形式 (9.1-2)，其留数等于 $\varpi(-\check\rho)$，其中 $\varpi$ 是投影 $\fh\to\fh/W$。由于映射 $\Op^{RS}_G(D)\to\Op_G(D^\times)$ 是单射（由命题 9.1.1），可知仿射空间 $\Op_{{}^L G}(D)$ 自然地实现为带正则奇点且留数为 $\varpi(-\check\rho)$ 的 opers 空间的子空间。\transnote{卡点：这句话把 $\Op_{{}^L G}(D)$ 等同于「留数为 $\varpi(-\check\rho)$ 的 opers」。背景：(9.1-2) 中 $\mathbf{v}(t)\in\fb[[t]]$ 正则，故 $\mathbf{v}(0)\in\fb$；取 $\mathbf{v}(t)=-\check\rho$ 即得正则 oper，其留数类为 $\varpi(-\check\rho)$。又因 $-\check\rho$ 与 $\check\rho$ 同属一个 $W$-轨道，$\varpi(\check\rho)=\varpi(-\check\rho)$，原书两种写法指同一纤维。结论：它就是 §9.3 中 (9.3-6) 取 $\check\lambda=0$、$w=1$ 的那一层，也是后文 $\check\lambda$-平移的起点。[来源：原书 p.255；p.266 (9.3-6)]}

% [Tutor Note: 这一段说明“正则 opers”$\Op_G(D)=\Op^0_G$ 恰是留数类为 $\varpi(-\check\rho)$ 的带正则奇点 opers 的特殊纤维：把 $\mathbf{v}(t)$ 用 $\check\rho(t)$ 共轭后，常数项被平移到 $p_{−1}−\check\rho$。9.2 节对一般的 $\check\lambda$ 重复同一手法，那里用 $(\check\lambda+\check\rho)(t)$ 代替 $\check\rho(t)$，于是留数平移成 $\varpi(-\check\lambda-\check\rho)$。]

\section{幂零 opers}

\begin{keypoints}
    \item \textbf{核心思想：}对 $\check\lambda$ 为整优余权的情形，用 $(\check\lambda+\check\rho)$-共轭把带正则奇点的 oper 化为另一实现，其留数带有“次级”不变量——$\mathfrak{n}/B$ 中的元素。
    \item \textbf{数学工具：}主分次、Coxeter 数 $h$、逐步规范变换消去系数、Deligne 扩张、$\check\lambda$-oper 条件。
    \item \textbf{关键结论：}命题 9.2.1：$\delta_{\check\lambda}$ 是同构；推论 9.2.3 给出 $\Op^{\mathrm{nilp}}_G$ 的坐标描述；命题 9.2.6 用 Deligne 扩张刻画幂零 opers。
    \item \textbf{记号提醒：}$\Op^{\mathrm{nilp},\check\lambda}_G$ 是余权 $\check\lambda$ 的幂零 opers；$\Op^{\check\lambda}_G$ 是留数映射为零（单值性平凡）的纤维；$\check\lambda=0$ 时简记 $\Op^{\mathrm{nilp}}_G$。
\end{keypoints}

现在我们引入带正则奇点 opers 的另一种实现，其留数形如 $\varpi(-\check\lambda-\check\rho)$，其中 $\check\lambda$ 是整优权（即满足 $\langle\alpha_i,\check\lambda\rangle\in\Z_+$），而 $\varpi(-\check\lambda-\check\rho)$ 是 $(-\check\lambda-\check\rho)\in\fh$ 在 $\fh/W$ 上的投影。这种实现对我们将很方便，因为它揭示了 oper 的一个“次级”留数不变量，它是 $\mathfrak{n}/B$ 中的元素（与上面定义的、$\fh/W$ 中的“初级”留数相对）。\transnote{卡点：$\check\lambda$ 到底是「权」还是「余权」。背景：原书自身混用——p.255 此处印 dominant integral weight，p.257 命题 9.2.1 印 coweight，p.263 又印 weight。由于 $\fh$ 上的不变内积给出 $\fh\cong\fh^*$，两种说法指同一个元素 $\check\lambda\in\fh$，数学上并无差别。结论：译文按各页页图分别照印「整优权／整优余权」，不强行统一，读者按 $\fh$ 中元素理解即可。[来源：原书 p.255、p.257、p.263 页图]}

为说明两种实现的差别，考虑 $\fg=\fsl_2$ 的情形。和以前一样，用 $\check\rho\mapsto 1$ 把 $\fh$ 与 $\C$ 等同。留数为 $-1-\check\lambda$（其中 $\check\lambda\in\Z_+$）的 $\fsl_2$-oper，是形如
\[
\partial_t + \begin{pmatrix} -\frac{1+\check\lambda}{2t}+a(t) & b(t)\\[2pt] \frac{1}{t} & \frac{1+\check\lambda}{2t}-a(t)\end{pmatrix},
\]
的算子的 $N[[t]]$-规范等价类，其中 $a(t),b(t)\in\C[[t]]$。用命题 9.1.1 证明中的规范变换，可以把这个算子化为典范形。具体地说，先作矩阵
\[
\begin{pmatrix} 1 & -\frac{1}{2t}\\[2pt] 0 & 1\end{pmatrix}
\begin{pmatrix} t^{-1/2} & 0\\ 0 & t^{1/2}\end{pmatrix}
\begin{pmatrix} 1 & \frac{1+\check\lambda}{2}-ta(t)\\[2pt] 0 & 1\end{pmatrix}
\]
的规范变换，得到典范形算子
\[
\partial_t + \begin{pmatrix} 0 & \frac{1}{t^2}\left(\frac{1+\check\lambda}{2}-ta(t)\right)^2-\frac{1}{4t^2}+\frac{1}{t}(b(t)+a(t))+a'(t)\\[3pt] 1 & 0\end{pmatrix}.
\]
另一方面，若作矩阵
\[
(\check\rho+\check\lambda)(t)^{-1} = \begin{pmatrix} t^{(1-\check\lambda)/2} & 0\\ 0 & t^{(1+\check\lambda)/2}\end{pmatrix}
\]
的规范变换，便得到算子
\[
\partial_t + \begin{pmatrix} a(t) & \dfrac{b(t)}{t^{\check\lambda+1}}\\[3pt] t^{\check\lambda} & -a(t)\end{pmatrix}.
\]
设 $b(t) = \sum_{n\ge 0}b_nt^n$。再作矩阵
\[
\begin{pmatrix} 1 & \sum_{n=0}^{\check\lambda-1}\dfrac{b_n}{-\check\lambda+n}t^{-\check\lambda+n}\\[3pt] 0 & 1\end{pmatrix}
\]
的规范变换，便得到形如
\[
\partial_t + \begin{pmatrix} \tilde a(t) & \dfrac{\tilde b(t)}{t}\\[3pt] t^{\check\lambda} & -\tilde a(t)\end{pmatrix}
\]
的算子，其中 $\tilde a(t),\tilde b(t)\in\C[[t]]$。

这是幂零 oper 的一个例子。它作为联络（不要与作为 oper 的留数混淆！）的留数是幂零矩阵
\[
\begin{pmatrix} 0 & \tilde b_0\\ 0 & 0\end{pmatrix},
\]
它在 $B$-共轭下良定义。由此可知，这个 oper 的单值性与矩阵
\[
\begin{pmatrix} 1 & 2\pi\ii\tilde b_0\\ 0 & 1\end{pmatrix}
\]
共轭。

% [Tutor Note: 请区分同一对象的两个“留数”：作为联络的留数（$\mathfrak{n}$ 中的幂零矩阵，在 B-共轭下良定义，决定单值性）与作为 oper 的留数（$p_{−1}+\mathbf v(0)$ 在 $\mathfrak{h}/W$ 中的类，见 9.1.2 节）。本书用同一词“residue”，译文保留“留数”，读者须按上下文分辨。]

于是，除典范形之外，我们得到留数为整的 opers 的另一种实现。这种实现的优点是有自然的留数映射，因而能更清楚地看出这些 opers 的单值性质。\transnote{卡点：residue 的两义，以及「留数为整」的所指。背景：同一算子有两套留数：§9.1 的 oper 留数是它在共轭类空间 $\fg/G=\fh/W$ 中的像（$\operatorname{res}$）；§9.2 的联络留数是 (9.2-1) 中的 $\mathbf{v}(0)\in\fn$（$\operatorname{Res}$），它决定单值性。「留数为整」说的是 $\check\lambda$ 整，故 (9.2-1) 中 $t^{\langle\alpha_i,\check\lambda\rangle}$ 是整数次幂。结论：这正是 §9.2 全部构造的出发点。[来源：原书 p.256（原文明确要求勿与 oper 的留数混淆）、p.262]}

本节按 [Frenkel and Gaitsgory 2006c] 说明如何把这一构造推广到其他单李代数。

\subsection{定义}

设 $\check\lambda\in\fh$ 是整优余权。考虑形如
\begin{equation}
\nabla = \partial_t + \sum_{i=1}^{\ell}t^{\langle\alpha_i,\check\lambda\rangle}\psi_i(t)f_i + \mathbf{v}(t) + \frac{v}{t},
\tag{9.2-1}
\end{equation}
的算子构成的空间，其中 $\psi_i(t)\in\C[[t]]$，$\psi_i(0)\ne 0$，$\mathbf{v}(t)\in\fb[[t]]$，$v\in\fn$。群 $B[[t]]$ 通过规范变换作用在这个空间上。由于 $B[[t]]$ 是群 $B((t))$ 的子群，而后者自由地作用在一切形如 (4.2-1) 的算子空间上（见 4.2.4 节），可知 $B[[t]]$ 的这个作用也是自由的。把一切形如 (9.2-1) 的算子空间商去 $B[[t]]$，所得空间称为余权 $\check\lambda$ 的\textbf{幂零 $G$-opers}空间，记作 $\Op^{\mathrm{nilp},\check\lambda}_G$。

我们有映射
\[
\beta_{\check\lambda}: \Op^{\mathrm{nilp},\check\lambda}_G\to\Op_G(D^\times),
\]
它把形如 (9.2-1) 的算子 $\nabla$ 的 $B[[t]]$-等价类映到其 $B((t))$-等价类。这个映射经过映射
\[
\delta_{\check\lambda}: \Op^{\mathrm{nilp},\check\lambda}_G\to\Op^{RS}_G(D)_{\varpi(-\check\lambda-\check\rho)}
\]
分解，其中 $\varpi(-\check\lambda-\check\rho)$ 是 $(-\check\lambda-\check\rho)\in\fh$ 在 $\fh/W$ 上的投影。$\delta_{\check\lambda}$ 构造如下。

给定形如 (9.2-1) 的算子 $\nabla$，考虑算子
\begin{equation}
\begin{gathered}
(\check\lambda+\check\rho)(t)\nabla(\check\lambda+\check\rho)(t)^{-1}
= \partial_t + t^{-1}\left(\sum_{i=1}^{\ell}\psi_i(t)f_i + \mathbf{w}(t)\right),\\
\mathbf{w}(t)\in(\check\lambda+\check\rho)(t)\fb[[t]](\check\lambda+\check\rho)(t)^{-1},\qquad \mathbf{w}(0) = -\check\lambda-\check\rho.
\end{gathered}
\tag{9.2-2}
\end{equation}
此外，在这个共轭下，$\nabla$ 的 $B[[t]]$-等价类被映到 $(\check\lambda+\check\rho)(t)\nabla(\check\lambda+\check\rho)(t)^{-1}$ 关于子群
\[
(\check\lambda+\check\rho)(t)B[[t]](\check\lambda+\check\rho)(t)^{-1}\subset B[[t]]
\]
的等价类。于是，所要的映射 $\delta_{\check\lambda}$ 把 $\nabla$ 的 $B[[t]]$-等价类映到 $(\check\lambda+\check\rho)(t)\nabla(\check\lambda+\check\rho)(t)^{-1}$ 的 $B[[t]]$-等价类。按构造，后者是带正则奇点、留数为 $\varpi(-\check\lambda-\check\rho)$ 的 oper。

\subsection{幂零 opers 与带正则奇点的 opers}

显然，映射 $\beta_{\check\lambda}$ 是 $\delta_{\check\lambda}$ 与包含映射
\[
\Op^{RS}_G(D)_{\varpi(-\check\lambda-\check\rho)}\hookrightarrow\Op_G(D^\times)
\]
的复合。

\begin{proposition}
对任意整优余权 $\check\lambda$，映射 $\delta_{\check\lambda}$ 是同构。
\end{proposition}

\begin{proof}
我们将证明：对任意 $\C$-代数 $R$，态射 $\delta_{\check\lambda}$ 在 $R$-点层面给出双射。这将蕴含 $\delta_{\check\lambda}$ 不仅是集合意义上的同构，而且作为概形态射也是同构。

于是需要证明：任何形如 (9.1-1)、留数为 $\varpi(-\check\lambda-\check\rho)$ 的算子，在任意 $R$ 上都可以化为 (9.2-2) 的形式；并且，若 (9.2-2) 形式的两个算子在 $g\in B(R[[t]])$ 的规范作用下共轭，则 $g\in(\check\lambda+\check\rho)(t)B(R[[t]])(\check\lambda+\check\rho)(t)^{-1}$。

采用记号
\[
\check\mu = \check\lambda+\check\rho = \sum_i\check\mu_i\check\omega_i.
\]
按对 $\check\lambda$ 的假设，对一切 $i$ 有 $\check\mu_i\ge 1$。

注意
\[
(\check\lambda+\check\rho)(t)B(R[[t]])(\check\lambda+\check\rho)(t)^{-1}
= H(R[[t]])\cdot(\check\lambda+\check\rho)(t)N(R[[t]])(\check\lambda+\check\rho)(t)^{-1}.
\]
存在唯一的 $H(R[[t]])$ 中元素，使 (9.1-1) 中的每个函数 $\psi_i(t)$ 都等于 1。因此只需证明：任何形如 (9.1-2)、留数为 $\varpi(-\check\lambda-\check\rho)$ 的算子 $\nabla$，在 $N(R[[t]])$ 的作用下都可化为
\begin{equation}
\nabla = \partial_t + t^{-1}(p_{-1}+\mathbf{v}(t)),\qquad
\mathbf{v}(t)\in\check\mu(t)\fb[[t]]\check\mu(t)^{-1},\quad \mathbf{v}(0)=-\check\mu,
\tag{9.2-3}
\end{equation}
的形式；并且，若这种形式的两个算子在 $g\in N(R[[t]])$ 的规范作用下共轭，则 $g\in\check\mu(t)N(R[[t]])\check\mu(t)^{-1}$。为简化记号，我们省略 $R$，把 $N(R[[t]])$ 简记作 $N[[t]]$，等等。

先考虑 $\check\lambda=0$，即对一切 $i$ 有 $\check\mu_i=1$ 的情形。我们用归纳法把形如 (9.1-2) 的算子化为 (9.2-3)。有分解
\[
\check\rho(t)\fb[[t]]\check\rho(t)^{-1} = \fh[[t]]\oplus\bigoplus_{j=1}^{h}t^j\fn_j[[t]],
\]
其中 $\fn = \bigoplus_{j>0}\fn_j$ 是 $\fn$ 关于主分次的分解，$h$ 是 Coxeter 数（见 4.2.4 节）。我们将假设 $h>1$，否则 $\fg=\fsl_2$，算子已在所要的形式中。

回忆按我们的假设，$\nabla$ 的留数等于 $\varpi(-\check\rho)$。因此可用常值子群 $N\subset N[[t]]$ 的规范作用把 $\nabla$ 化为
\begin{equation}
\partial_t + \frac{1}{t}(p_{-1}-\check\rho) + \mathbf{v}(t),
\tag{9.2-4}
\end{equation}
其中 $\mathbf{v}(t)\in\fb[[t]]$。按主分次把 $\mathbf{v}(t)$ 分解为 $\mathbf{v}(t)=\sum_{j=0}^{h}\mathbf{v}_j(t)$，并作
\[
g = \exp\left(-\frac{1}{h-1}t\mathbf{v}_h(0)\right)\in N[[t]]
\]
的规范变换。新算子中主次数 $h$ 项（即 $\mathbf{v}_h$）的常数系数为
\[
\mathbf{v}_h(0) + \frac{1}{h-1}[-\mathbf{v}_h(0),-\check\rho] + \frac{1}{h-1}\mathbf{v}_h(0) = 0
\]
（第二项来自 $-g\frac{\check\rho}{t}g^{-1}$，最后一项来自 $g\partial_t(g^{-1})$）。于是新算子的 $\mathbf{v}_{\mathrm{new},h}(t)\in t\fn_h[[t]]$。

若 $h=2$，则已完毕。否则再作 $\exp(-\frac{1}{h-2}t\mathbf{v}_{h-1}(0))$ 的规范变换（这里及下文用 $\mathbf{v}_i(t)$ 表示最近一次规范变换后所得算子的分量）。同样得到新算子的 $\mathbf{v}_i(t)\in t\fn_i[[t]]$，$i=h-1,h$。如此继续，得到形如 (9.2-4) 的算子，其中 $\mathbf{v}_1(t)\in\fn_1[[t]]$，且当 $j=2,\dots,h$ 时 $\mathbf{v}_j(t)\in t\fn_j[[t]]$。（容易看出，每次规范变换都保持前几步得到的消没条件。）

接着，我们希望得到形如 (9.2-4) 的算子，其中
\[
\mathbf{v}_1(t)\in\fn_1[[t]],\qquad \mathbf{v}_2(t)\in t\C[[t]],\qquad
\mathbf{v}_j(t)\in t^2\fn_j[[t]]\quad (j=3,\dots,h).
\]
这按与上面相同的方式完成：依次作
\[
\exp\left(-\frac{1}{j-2}t^2\mathbf{v}'_j(0)\right),\qquad j=h,h-1,\dots,3,
\]
的规范变换，其中 $\mathbf{v}'_j(0)$ 表示前一步所得算子第 $j$ 个分量的 $t$-系数。显然，这一程序对任意 $\C$-代数 $R$ 都有效。

如此进行下去，把算子化为
\begin{equation}
\partial_t + \frac{1}{t}(p_{-1}+\mathbf{v}(t)),\qquad
\mathbf{v}(t)\in\check\rho(t)\fb[[t]]\check\rho(t)^{-1},\quad \mathbf{v}(0)=-\check\rho,
\tag{9.2-5}
\end{equation}
即 $\check\mu=\check\rho$ 时的 (9.2-3)。下面证明：若 $\nabla$ 与 $g\nabla g^{-1}$ 都有形式 (9.2-5) 且 $g\in N[[t]]$，则 $g$ 必属于子群 $\check\rho(t)N[[t]]\check\rho(t)^{-1}$（这等价于说：初始算子 $\nabla$ 的 $N[[t]]$-等价类与一切形如 (9.2-5)（其中 $\check\mu=\check\rho$）的算子集合之交，是一个 $\check\rho(t)N[[t]]\check\rho(t)^{-1}$-等价类）。

由于常值子群 $N\subset N[[t]]$ 中任何非恒等元素的规范作用都会改变联络的留数，不妨设 $g$ 属于 $N[[t]]$ 的第一同余子群 $N^{(1)}[[t]]$。由于指数映射 $\exp: t\fn[[t]]\to N^{(1)}[[t]]$ 是同构，可以把 $N^{(1)}[[t]]$ 中任一元素唯一地写成形如 $\exp(c_{\alpha,n}t^n)$（$c_{\alpha,n}\in\fn_\alpha$）的元素之积，顺序任意。特别地，可以把 $N^{(1)}[[t]]$ 中任一元素写成两个元素之积，其一属于子群 $\check\rho(t)N[[t]]\check\rho(t)^{-1}$，另一个是形如 $\exp(c_{\alpha,n}t^n)$（$\alpha\in\Delta_+$，$1\le n<\langle\alpha,\check\rho\rangle$）的元素之积。但把算子化为 (9.2-5) 所用的正是第二类元素的规范变换，且它们由这一程序唯一确定。因此对 (9.2-5) 施加其中任一元素都会改变其形式。所以 $N^{(1)}[[t]]$ 中保持形式 (9.2-5) 的变换构成的子群恰是子群 $\check\mu(t)N[[t]]\check\mu(t)^{-1}$。

这就完成了 $\check\lambda=0$、$\check\mu=\check\rho$ 情形的证明。

对一般的 $\check\mu$，用归纳法证明命题的断言。先注意，与上面相同的论证说明：任何形如 (9.1-1)、留数为 $\varpi(\check\mu)$ 的算子都可化为
\[
\partial_t + t^{-1}(p_{-1}+\mathbf{v}(t)),\qquad
\mathbf{v}(t)\in\check\rho(t)\fb[[t]]\check\rho(t)^{-1},\quad \mathbf{v}(0)=-\check\mu,
\]
其中对任意 $\check\mu=\check\lambda+\check\rho$（$\check\lambda$ 为整优余权）都成立。\transnote{卡点：这里的 $\varpi(\check\mu)$ 少一个负号。背景：原书 p.260 顶行印作 residue $\varpi(\check\mu)$，但同页 (9.2-3) 明确 $\mathbf{v}(0)=-\check\mu$，p.261 末又写作 $\varpi(-\check\mu)$。结论：按上下文应为 $\varpi(-\check\mu)$，这是原书笔误；译文照印原文，未改动公式。[来源：原书 p.260 与 p.261 页图]}我们证明这个算子还可进一步化为
\begin{equation}
\partial_t + t^{-1}(p_{-1}+\mathbf{v}(t)),\qquad
\mathbf{v}(t)\in\check\mu'(t)\fb[[t]]\check\mu'(t)^{-1},\quad \mathbf{v}(0)=-\check\mu,
\tag{9.2-6}
\end{equation}
其中 $\check\mu'$ 是满足
\begin{equation}
1\le\langle\alpha_i,\check\mu'\rangle\le\langle\alpha_i,\check\mu\rangle,\qquad i=1,\dots,\ell
\tag{9.2-7}
\end{equation}
的任意整优余权。

% [Tutor Note: 从 p.260 顶行起，原书在此处把留数印作 $\varpi(\check\mu)$（无负号），但同一证明在 p.261 末尾又写“留数为 $\varpi(-\check\mu)$ 的算子”，且目标形式 (9.2-3) 的 $\mathbf{v}(0)=-\check\mu$ 给出留数类 $\varpi(-\check\mu)$。故 p.260 顶行的 $\varpi(\check\mu)$ 疑为 $\varpi(-\check\mu)$ 之误；译文保留原文并在此说明。]

% [Tutor Note: 归纳变量是 $\check\mu'$（不是 $\check\mu$），条件 (9.2-7) 的两边分别由 $\check\mu$ 与 $\check\mu'$ 控制：$1\le\langle\alpha_i,\check\mu'\rangle\le\langle\alpha_i,\check\mu\rangle$。归纳步每次把某个 $\check\mu'_i$ 加 1，直到 $\check\mu'=\check\mu$；每一步用 $t^{\check\mu'_i}$ 阶规范变换把 $\alpha_i$ 方向的系数推到更高一阶。读这段时容易把归纳对象误当成 $\check\mu$。]
已知当 $\check\mu'=\check\rho$ 时这成立。设对某个 $\check\mu'$ 已证；下面证明对余权 $\check\mu'+\check\omega_i$ 也成立（假定它仍满足条件 (9.2-7)）。\transnote{卡点：原书这一句印作「对余权 $\check\mu+\check\omega_i$ 也成立」。背景：同段紧接着说「希望把它化为同一形式，但把 $\check\mu'$ 换成 $\check\mu'+\check\omega_i$」，且条件 (9.2-7) 的变量是 $\check\mu'$；归纳变量只能是 $\check\mu'$，写成 $\check\mu$ 会与已固定的 $\check\mu=\check\lambda+\check\rho$ 冲突。结论：原书此处漏一撇，应为 $\check\mu'+\check\omega_i$；译文按此处理。[来源：原书 p.260 同段两句页图]}于是有形式 (9.2-6) 的算子，希望把它化为同一形式，但把 $\check\mu'$ 换成 $\check\mu'+\check\omega_i$。利用根分解 $\fn = \bigoplus_{\alpha\in\Delta_+}\fn_\alpha$，把 $\mathbf{v}(t)$ 写成 $\mathbf{v}(t) = \sum_{\alpha\in\Delta_+}\mathbf{v}_\alpha(t)$，其中 $\mathbf{v}_\alpha(t)\in\fn_\alpha$。有
\begin{equation}
\mathbf{v}_\alpha(t) = \sum_{n\ge\langle\alpha,\check\mu'\rangle}\mathbf{v}_{\alpha,n}t^n.
\tag{9.2-8}
\end{equation}
我们逐步归纳。第一步考虑 $\alpha$ 为单根时的分量 $\mathbf{v}_\alpha(t)$（即属于 $\fn_1$ 的那些）。对一切这样的 $\alpha$ 都有 $\langle\alpha,\check\mu'+\check\omega_i\rangle = \langle\alpha,\check\mu'\rangle$，例外是 $\alpha=\alpha_i$，此时 $\langle\alpha,\check\mu'+\check\omega_i\rangle = \langle\alpha,\check\mu'\rangle+1$。

记 $\check\mu_i = \langle\alpha_i,\check\mu\rangle$、$\check\mu'_i = \langle\alpha_i,\check\mu'\rangle$。则对 $\check\mu'+\check\omega_i$ 的条件 (9.2-7) 变成
\begin{equation}
2\le\check\mu'_i+1\le\check\mu_i.
\tag{9.2-9}
\end{equation}
作
\[
\exp\left(\frac{1}{\check\mu'_i-\check\mu_i}t^{\check\mu'_i}\mathbf{v}_{\alpha_i,\check\mu'_i}\right)
\]
的规范变换（注意由条件 (9.2-9)，分母非零）。结果得到形如 (9.2-6) 的算子，其中 $\mathbf{v}_{\alpha_i}(t)\in t^{\check\mu'_i+1}\C[[t]]$。同时，对其他任何 $\alpha\in\Delta_+$，条件 (9.2-8) 不会被破坏，因为按对 $\check\mu'$ 的假设有 $\langle\alpha,\check\mu'\rangle\ge\langle\alpha,\check\rho\rangle$。

接着处理高度为 2 的根（即属于 $\fn_2$ 的那些）。对满足 $\langle\alpha,\check\mu'+\check\omega_i\rangle>\langle\alpha,\check\mu'\rangle$ 的 $\alpha$，依次作
\[
\exp\left(\frac{1}{\langle\alpha,\check\mu'\rangle-\langle\alpha,\check\mu\rangle}t^{\langle\alpha,\check\mu'\rangle}\mathbf{v}_{\alpha,\langle\alpha,\check\mu'\rangle}\right)
\]
的规范变换。沿主分次的所有值继续这一过程，得到这样的算子：对一切满足 $\langle\alpha,\check\mu'\rangle = \langle\alpha,\check\mu'+\check\omega_i\rangle$ 的 $\alpha$ 都满足条件 (9.2-8)；而对一切满足 $\langle\alpha,\check\mu'+\check\omega_i\rangle>\langle\alpha,\check\mu'\rangle$ 的 $\alpha$ 满足
\[
\mathbf{v}_\alpha(t) = \sum_{n\ge\langle\alpha,\check\mu'\rangle+1}\mathbf{v}_{\alpha,n}t^n.
\]
下一步用类似方式得到满足更强条件的算子：除上述之外，对满足 $\langle\alpha,\check\mu'+\check\omega_i\rangle>\langle\alpha,\check\mu'\rangle+1$ 的 $\alpha$ 还有
\[
\mathbf{v}_\alpha(t) = \sum_{n\ge\langle\alpha,\check\mu'\rangle+2}\mathbf{v}_{\alpha,n}t^n.
\]
如此继续，最终得到对一切 $\alpha\in\Delta_+$ 都有
\[
\mathbf{v}_\alpha(t) = \sum_{n\ge\langle\alpha,\check\mu'+\check\omega_i\rangle}\mathbf{v}_{\alpha,n}t^n
\]
的算子。换言之，它具有形式 (9.2-3)，只是把 $\check\mu'$ 换成了 $\check\mu'+\check\omega_i$。这就完成了归纳步。\transnote{卡点：此处引的编号 (9.2-3)。背景：把 $\check\mu'$ 换成 $\check\mu'+\check\omega_i$ 后得到的是 (9.2-6) 的形状（其中 $\check\mu'$ 处代之以 $\check\mu'+\check\omega_i$）；(9.2-3) 是 $\check\mu=\check\lambda+\check\rho$ 的特殊情形，不含可变的 $\check\mu'$。结论：原书此处编号笔误，应为 (9.2-6)；译文照印不改。[来源：原书 p.261 与 p.260 (9.2-6)]}

于是得到：任何形如 (9.1-2)、留数为 $\varpi(-\check\mu)$ 的算子都可化为 (9.2-3)（显然，这一构造对任意 $\C$-代数 $R$ 都有效）。最后，与 $\check\mu=\check\rho$ 情形相同可证：若这种形式的两个算子在 $g\in N[[t]]$ 的规范作用下共轭，则 $g\in\check\mu(t)N[[t]]\check\mu(t)^{-1}$。
\end{proof}

由命题 9.1.1 与命题 9.2.1，映射 $\beta_{\check\lambda}:\Op^{\mathrm{nilp},\check\lambda}_G\to\Op_G(D^\times)$ 是单射。特别地，利用命题 9.1.1 可得幂零 opers 的典范代表。下文中我们不再区分 $\Op^{\mathrm{nilp},\check\lambda}_G$ 与它在 $\Op_G(D^\times)$ 中的像。

\subsection{留数}

设 $\mathcal{F}_B$ 是圆盘 $D$ 上承载 oper $\chi\in\Op^{\mathrm{nilp},\check\lambda}_G$ 的 $B$-丛，$\mathcal{F}_{B,0}$ 是它在 $0\in D$ 处的纤维。与 4.2.4 节相同可证：$D$ 上所有 $\check\lambda$-幂零 opers 底下的 $B$-丛 $\mathcal{F}_B$ 被典范等同。把 $\mathfrak{n}$ 的 $\mathcal{F}_{B,0}$-扭变记作 $\mathfrak{n}_{\mathcal{F}_{B,0}}$。则对每个形如 (9.2-1) 的算子，$v$ 是 $\mathfrak{n}_{\mathcal{F}_{B,0}}$ 中良定义的向量。因此得到态射
\[
\operatorname{Res}_{\mathcal{F},\check\lambda}: \Op^{\mathrm{nilp},\check\lambda}_G\to\mathfrak{n}_{\mathcal{F}_{B,0}},\qquad \chi\mapsto v.
\]
对 $\mathcal{F}_{B,0}$ 的每个平凡化选取，得到等同 $\mathfrak{n}_{\mathcal{F}_{B,0}}\simeq\mathfrak{n}$。平凡化的更换导致 $B$ 在 $\mathfrak{n}$ 上的伴随作用。因此有典范映射 $\mathfrak{n}_{\mathcal{F}_{B,0}}\to\mathfrak{n}/B$。它与 $\operatorname{Res}_{\mathcal{F},\check\lambda}$ 的复合是态射
\[
\operatorname{Res}_{\check\lambda}: \Op^{\mathrm{nilp},\check\lambda}_G\to\mathfrak{n}/B.
\]
它把形如 (9.2-1) 的算子映到 $v$ 在 $\mathfrak{n}/B$ 中的投影。我们称之为\textbf{留数态射}。注意，任何幂零 $G$-oper 作为 $D^\times$ 上的联络都有单值算子，即 $G$ 中的一个共轭类。下述断言是显然的。

\begin{lemma}
任何余权为 $\check\lambda$ 的幂零 $G$-oper $\chi$ 都有幂幺单值性。其单值性平凡当且仅当 $\operatorname{Res}_{\check\lambda}(\chi)=0$。
\end{lemma}

把 $\Op^{\mathrm{nilp},\check\lambda}_G$ 中在映射 $\operatorname{Res}_{\check\lambda}$ 下 $0\in\mathfrak{n}/B$ 的原像记作 $\Op^{\check\lambda}_G$。等价地，它是形如
\begin{equation}
\nabla = \partial_t + \sum_{i=1}^{\ell}t^{\langle\alpha_i,\check\lambda\rangle}\psi_i(t)f_i + \mathbf{v}(t)
\tag{9.2-10}
\end{equation}
的算子的 $B[[t]]$-等价类构成的空间，其中 $\psi_i(t)\in\C[[t]]$，$\psi_i(0)\ne 0$，$\mathbf{v}(t)\in\fb[[t]]$。这个空间最初由 V. Drinfeld 引入（未发表）。

上述引理与命题 9.2.1 蕴含：$\Op^{\check\lambda}_G$ 是 $\Op^{RS}_G(D)_{\varpi(-\check\lambda-\check\rho)}$ 中对应单值性平凡的那些带正则奇点 opers 的子流形。

更详细地考虑 $\check\lambda=0$ 的情形。我们把 $\Op^{\mathrm{nilp},0}_G$ 简记作 $\Op^{\mathrm{nilp}}_G$，并称其点为\textbf{幂零 $G$-opers}。它们是形如
\begin{equation}
\partial_t + \sum_{i=1}^{\ell}\psi_i(t)f_i + \mathbf{v}(t) + \frac{v}{t},
\tag{9.2-11}
\end{equation}
的算子的 $B[[t]]$-等价类（其中 $\psi_i(t)\in\C[[t]]$，$\psi_i(0)\ne 0$，$\mathbf{v}(t)\in\fb[[t]]$，$v\in\fn$），或等价地，是形如
\begin{equation}
\partial_t + p_{-1} + \mathbf{v}(t) + \frac{v}{t}
\tag{9.2-12}
\end{equation}
的算子的 $N[[t]]$-等价类，其中 $\mathbf{v}(t)\in\fb[[t]]$，$v\in\fn$。留数态射（此时简记作 $\operatorname{Res}$）把这样的算子映到 $v$ 在 $\mathfrak{n}/B$ 中的像。注意 $\mathfrak{n}/B$ 中 $0$ 上的纤维恰是圆盘 $D$ 上正则 opers 构成的子集 $\Op^0_G = \Op_G(D)$。

命题 9.1.1 把 $\Op^{\mathrm{nilp}}_G$ 与 9.1.1 节定义的映射 $\operatorname{res}:\Op^{RS}_G(D)\to\fh/W$ 下 $\varpi(-\check\rho)$ 的原像等同起来。利用命题 9.2.1 与命题 9.1.1，得到 $\Op^{\mathrm{nilp}}_G$ 的如下描述。\transnote{卡点：推论 9.2.3 的上同调次数下界比推论 9.1.2 整体左移一位。背景：9.1.2 给 $\mathcal{P}roj_{\ge-2}\times\bigoplus_{j>1}\omega^{\otimes(d_j+1)}_{\ge-(d_j+1)}$，9.2.3 给 $\mathcal{P}roj_{\ge-1}\times\bigoplus_{j\ge2}\omega^{\otimes(d_j+1)}_{\ge-d_j}$。结论：差别来自常数项的自由度——(9.1-5) 把 $p_{-1}$ 放在常数项（阶 $0$），(9.2-3) 把它并入 $t^{-1}$ 的系数（阶 $-1$），于是各阶极点允许高一阶、求和指标也从 $j>1$ 改为 $j\ge2$。[来源：原书 p.262 对 p.254]}

\begin{corollary}
存在同构
\begin{equation}
\Op^{\mathrm{nilp}}_G\simeq \mathcal{P}roj_{\ge-1}\times\bigoplus_{j\ge 2}\omega^{\otimes(d_j+1)}_{\ge-d_j}.
\tag{9.2-13}
\end{equation}
\end{corollary}

% [Tutor Note: 注意推论 9.2.3 中的下界从推论 9.1.2 的 $\ge-2$、$\ge-(d_j+1)$ 各上移了一位，变为 $\mathcal{P}roj_{\ge-1}$ 与 $\omega_{\ge-d_j}$：这是因为幂零 opers 只保留留数满足 $\varpi(-\check\rho)$ 的特殊纤维（少了一个常数自由度）。核对下标时以此差一为检查点。]

\subsection{带正则奇点 opers 的更一般形式}

命题 9.2.1 证明中所用的论证，使我们能构造带正则奇点 opers 的更一般的代表。

考虑空间 $\Op^{RS}(D)_{\varpi(-\check\mu)}$，其中 $\check\mu$ 是 $\fh$ 中任意元素。$\Op^{RS}(D)_{\varpi(-\check\mu)}$ 中的元素等同于形如
\begin{equation}
\partial_t + \frac{1}{t}(p_{-1}-\check\mu) + \mathbf{v}(t),\qquad \mathbf{v}(t)\in\fb[[t]]
\tag{9.2-14}
\end{equation}
的算子的 $N^{(1)}[[t]]$-等价类。

设 $\check\lambda'$ 是满足
\begin{equation}
\langle\alpha,\check\lambda'\rangle<\langle\alpha,\check\mu\rangle,\qquad \text{若 } \langle\alpha,\check\mu\rangle\in\Z_{>0},\ \alpha\in\Delta_+
\tag{9.2-15}
\end{equation}
的整优余权，其中 $\Z_{>0}$ 是正整数集。

令 $\check\mu' = \check\lambda'+\check\rho$。把形如 (9.2-6) 的算子在群 $\check\mu'(t)N[[t]]\check\mu'(t)^{-1}$ 的规范作用下的规范等价类空间记作 $\Op(D)^{\check\mu'}_{\check\mu}$。则我们有自然映射
\[
\delta^{\check\mu'}_{\check\mu}: \Op(D)^{\check\mu'}_{\check\mu}\to\Op^{RS}(D)_{\varpi(-\check\mu)},
\]
它把这些等价类映到 $N[[t]]$-等价类。逐字照搬命题 9.2.1 证明中的论证，得到如下结果。

\begin{proposition}
若 $\check\lambda'$ 满足条件 (9.2-15)，且 $\check\mu'=\check\lambda'+\check\rho$，则映射 $\delta^{\check\mu'}_{\check\mu}$ 是同构。
\end{proposition}

特别地，设对一切 $\alpha\in\Delta_+$ 都有 $\langle\alpha,\check\mu\rangle\notin\Z_{>0}$。注意 $\fh/W$ 中任一点都可表示为 $\varpi(-\check\mu)$，其中 $\check\mu$ 满足这一条件。此时命题 9.2.4 的断言对任意整优权 $\check\lambda'$ 都成立。令 $\check\lambda' = n\check\rho$，$n\in\Z_+$。则 $\check\mu' = (n+1)\check\rho$，且 $\check\mu'(t)N[[t]]\check\mu'(t)^{-1}$ 是 $N[[t]]$ 的第 $(n+1)$ 个同余子群。取极限 $n\to\infty$，得到：

\begin{corollary}
设 $\check\mu$ 满足：对一切 $\alpha\in\Delta_+$ 都有 $\langle\alpha,\check\mu\rangle\notin\Z_{>0}$。则任何带正则奇点、留数为 $\varpi(-\check\mu)$ 的 oper 都可以唯一地写成
\[
\partial_t + \frac{1}{t}(p_{-1}-\check\mu) + \mathbf{v}(t),\qquad \mathbf{v}(t)\in\fh[[t]],
\]
或等价地，写成
\begin{equation}
\partial_t + p_{-1} - \frac{\check\mu+\check\rho}{t} + \mathbf{v}(t),\qquad \mathbf{v}(t)\in\fh[[t]]
\tag{9.2-16}
\end{equation}
的形式。于是，$\Op^{RS}(D)_{\varpi(-\check\mu)}$ 与形如 (9.2-16) 的算子空间之间存在双射。
\end{corollary}

\subsection{用 Deligne 扩张刻画幂零 opers}

回忆：给定去心圆盘 $D^\times$ 上带联络 $\nabla$ 的 $G$-丛 $\mathcal{F}$，若其单值性幂幺，则由 P. Deligne [1970] 定义的 $\mathcal{F}$ 到圆盘 $D$ 上 $G$-丛 $\overline{\mathcal{F}}$ 的典范扩张，使得联络有正则奇点且其余为幂零。考虑 $D^\times$ 上的 $G$-oper $(\mathcal{F},\mathcal{F}_B,\nabla)$。设它有幂幺单值性。$D^\times$ 上任何单值性幂幺的带联络 $G$-丛都有到 $D$ 上带正则奇点、幂零留数联络的丛的典范扩张，称为 Deligne 扩张。于是 $D^\times$ 上的丛 $\mathcal{F}$ 有到 $D$ 的 Deligne 扩张，记作 $\overline{\mathcal{F}}$。

由于旗簇是 紧合的，$\mathcal{F}$ 在 $D^\times$ 上的 $B$-约化 $\mathcal{F}_B$ 扩张为 $\overline{\mathcal{F}}$ 在 $D$ 上的 $B$-约化，记作 $\overline{\mathcal{F}}_B$。我们把 $(\overline{\mathcal{F}},\overline{\mathcal{F}}_B,\nabla)$ 称为 oper $(\mathcal{F},\mathcal{F}_B,\nabla)$ 到 $D$ 的 Deligne 扩张。按构造，联络 $\nabla$ 有正则奇点，其留数是
\[
\fg_{\overline{\mathcal{F}}_{B,0}} = \overline{\mathcal{F}}_B\times_{\overline{\mathcal{F}}_B}\fg
\]
中的元素，其中 $\overline{\mathcal{F}}_{B,0}$ 是 $\overline{\mathcal{F}}_B$ 在原点的纤维。
% [Tutor Note: 此处印作 $\overline{\mathcal{F}}_B\times_{\overline{\mathcal{F}}_B}\mathfrak{g}$，第二个下标疑为 $B$ 之误：与 $B$-丛相伴的 $\mathfrak{g}$-丛应写作 $\overline{\mathcal{F}}_B\times_B\mathfrak{g}$（$B$ 通过伴随作用）。两种写法给出同一纤维数值，但记号上应取 $B$。译文保留原文并登记。]

回忆 $H$-torsor $\mathbf{O}$，它是 $\mathfrak{n}^\perp/\fb\subset\fg/\fb$ 中的开稠子集。设有 $D$ 上取值于 $\mathfrak{n}^\perp/\fb$ 的函数 $f$，它在 $D^\times$ 上取值于 $\mathbf{O}$。给定整余权 $\check\lambda$，我们说 $f$ 在原点以 $\check\lambda$ 阶消没，如果相应的映射 $D^\times\to\mathbf{O}$ 可以作为余特征标 $\check\lambda:\C^\times\to\mathbf{O}$ 在某个嵌入 $D^\times\hookrightarrow\C^\times$ 与某个等同 $\mathbf{O}\simeq H$ 下的限制而得到。注意，若把取值于 $\mathfrak{n}^\perp/\fb$ 的函数换成 $\mathfrak{n}^\perp/\fb$ 关于 $D$ 上某个 $H$-丛的扭变的一个截面，这个定义仍然有意义。

% [Tutor Note: $\mathfrak{n}^\perp$ 是 $\mathfrak{b}$ 的正交补（相对于不变内积），于是 $\mathfrak{n}^\perp/\mathfrak{b}\subset\mathfrak{g}/\mathfrak{b}$ 是开稠的 $H$-torsor。这里“以 $\check\lambda$ 阶消没”的阶数是逐分量用 $\langle\nu,\check\lambda\rangle$ 度量的，与 9.3.3 节的 $\check\lambda$-次数是同一约定；下文 (9.3-14) 的计数正是按这一约定。]

我们说 $(\overline{\mathcal{F}},\overline{\mathcal{F}}_B,\nabla)$ 在原点满足 $\check\lambda$-oper 条件，如果 $D$ 上取值于 $(\mathfrak{n}^\perp/\fb)_{\overline{\mathcal{F}}_B}$（且在 $D^\times$ 上取值于 $\mathbf{O}_{\overline{\mathcal{F}}_B}\subset(\fg/\fb)_{\overline{\mathcal{F}}_B}$）的一次形式 $\nabla/\overline{\mathcal{F}}_B$ 在原点以 $\check\lambda$ 阶消没。下述命题直接由定义得出。

\begin{proposition}
设 $\check\lambda$ 是整优余权。$D^\times$ 上的 $G$-oper 是余权 $\check\lambda$ 的幂零 $G$-oper，当且仅当它有幂幺单值性且其 Deligne 扩张在原点满足 $\check\lambda$-oper 条件。
\end{proposition}

\section{带正则奇点的 Miura opers}

\begin{keypoints}
    \item \textbf{核心思想：}把 8.2.1 节的 Miura 变换限制到带正则奇点的 Cartan 联络，研究其纤维；幂零 oper 上方的纤维恰是留数的 Springer 纤维。
    \item \textbf{数学工具：}H-丛 $\Omega^{\check\rho}$ 上的联络、Miura 变换 $\mu$、$\check\lambda$-次数、Bruhat 分解与泛位置、Deligne 扩张。
    \item \textbf{关键结论：}式 (9.3-6) 给出 $\mu$ 的纤维的 $W$-分层；定理 9.3.6：$\tilde\mu_{-w(\check\lambda+\check\rho)+\check\rho}$ 是同构，且 $\mathrm{MOp}^{\check\lambda}_G\simeq\Op^{\mathrm{nilp},\check\lambda}_G\times_{\mathfrak{n}/B}\tilde{\mathfrak{n}}/B$。
    \item \textbf{记号提醒：}$\mu$ 是 Miura 变换，$\mu^{\mathrm{RS}}$ 是它在正则奇点情形的限制；$\tilde{\mathfrak{n}}$ 是 Springer 簇在 $\mathfrak{n}$ 上的限制；$\mathbf{O}$ 是 $\mathfrak{h}$ 的 $H$-torsor。
\end{keypoints}

8.2.1 节引入了 Miura opers 与 Cartan 联络的空间，以及到 opers 空间的 Miura 变换。本节考虑带正则奇点的 Miura opers 与 Cartan 联络，以及 Miura 变换在带正则奇点 Cartan 联络空间上的限制。我们证明：幂零 oper $\chi$ 上 Miura 变换的纤维同构于 $\chi$ 的留数的 Springer 纤维。本节结果由 [Frenkel and Gaitsgory 2006c] 得到。

\subsection{联络与带正则奇点的 opers}

8.2.1 节定义了圆盘 $D=\Spec\C[[t]]$ 上 $H$-丛 $\Omega^{\check\rho}$ 的联络空间 $\operatorname{Conn}(\Omega^{\check\rho})_D$。现在考虑该丛在去心圆盘 $D^\times = \Spec\C((t))$ 上的联络空间 $\operatorname{Conn}(\Omega^{\check\rho})_{D^\times}$。$\operatorname{Conn}(\Omega^{\check\rho})_{D^\times}$ 的元素由形如
\begin{equation}
\overline{\nabla} = \partial_t + \mathbf{u}(t),\qquad \mathbf{u}(t)\in\fh((t))
\tag{9.3-1}
\end{equation}
的算子代表；在坐标变换下它们按公式 (8.2-2) 变换。

现在令 $\operatorname{Conn}(\Omega^{\check\rho})^{\mathrm{RS}}_D$ 为 $D$ 上 $H$-丛 $\Omega^{\check\rho}$ 的一切带\textbf{正则奇点}的联络空间，即形如
\begin{equation}
\overline{\nabla} = \partial_t + \frac{1}{t}\mathbf{u}(t),\qquad \mathbf{u}(t)\in\fh[[t]]
\tag{9.3-2}
\end{equation}
的联络构成的空间。

我们有映射
\[
\operatorname{res}_{\fh}: \operatorname{Conn}(\Omega^{\check\rho})^{\mathrm{RS}}_D\to\fh,
\]
它把联络映到其留数 $\mathbf{u}(0)\in\fh$。对 $\check\lambda\in\fh$，把留数为 $\check\lambda$ 的联络子空间记作 $\operatorname{Conn}(\Omega^{\check\rho})^{\check\lambda}_D$。

回忆 Miura 变换 (8.3-3)
\begin{equation}
\mu: \operatorname{Conn}(\Omega^{\check\rho})_{D^\times}\to\Op_G(D^\times),
\tag{9.3-3}
\end{equation}
其定义为：给定形如 (9.3-1) 的算子 $\overline{\nabla}$，相应的 $G$-oper 定义为算子 $\overline{\nabla}+p_{-1}$ 的 $N((t))$-规范等价类。

设 $\overline{\nabla}$ 是由公式 (9.3-2) 给出的 $\operatorname{Conn}(\Omega^{\check\rho})^{\mathrm{RS}}_D$ 中元素。则 $\mu(\overline{\nabla})$ 按定义是算子
\[
\nabla = \partial_t + p_{-1} + t^{-1}\mathbf{u}(t)
\]
的 $N((t))$-等价类。用 $\check\rho(t)$ 作规范变换，把它与算子
\[
\partial_t + t^{-1}(p_{-1}-\check\rho+\mathbf{u}(t))
\]
的 $N((t))$-等价类等同，后者是带正则奇点、留数为 $\varpi(-\check\rho+\mathbf{u}(0))$ 的 oper。

于是得到态射
\begin{equation}
\mu^{\mathrm{RS}}: \operatorname{Conn}(\Omega^{\check\rho})^{\mathrm{RS}}_D\to\Op^{RS}_G(D),
\tag{9.3-4}
\end{equation}
以及交换图表
\begin{equation}
\begin{CD}
\operatorname{Conn}(\Omega^{\check\rho})^{\mathrm{RS}}_D @>\mu^{\mathrm{RS}}>> \Op^{RS}_G(D)\\
@V{\operatorname{res}_{\fh}}VV @VV{\operatorname{res}}V\\
\fh @>>> \fh/W
\end{CD}
\tag{9.3-5}
\end{equation}
其中下水平映射是映射 $\check\lambda\mapsto\check\lambda-\check\rho$ 与投影 $\varpi:\fh\to\fh/W$ 的复合。

容易看出，$\Op^{RS}_G(D)$ 中任一元素在 Miura 变换 (9.3-3) 下的原像属于 $\operatorname{Conn}(\Omega^{\check\rho})^{\mathrm{RS}}_D\subset\operatorname{Conn}(\Omega^{\check\rho})_{D^\times}$。换言之，带非正则奇点的 Cartan 联络所对应的 oper 必然有非正则奇点。因此，$\mu^{\mathrm{RS}}$ 是 $\mu$ 在 $\Op^{RS}_G(D)$ 上的限制。上述交换图表于是蕴含：对任意 $\check\lambda\in\fh$，有同构\footnote{本章中我们始终考虑 $\mu$ 的集合论原像。相应的概形论原像可能非既约（例子见 10.4.6 节）。}
\begin{equation}
\mu^{-1}\left(\Op^{RS}_G(D)_{\varpi(-\check\lambda-\check\rho)}\right)
= \bigsqcup_{w\in W}\operatorname{Conn}(\Omega^{\check\rho})^{\check\rho-w(\check\lambda+\check\rho)}_D.
\tag{9.3-6}
\end{equation}

此外，推论 9.2.5 蕴含：

\begin{proposition}
设 $\check\lambda\in\fh$ 满足：对一切 $\alpha\in\Delta_+$ 都有 $\langle\alpha,\check\lambda\rangle\notin\Z_+$。则映射 $\mu^{\mathrm{RS}}$ 在 $\operatorname{Conn}(\Omega^{\check\rho})^{-\check\lambda}_D$ 上的限制是同构
\[
\operatorname{Conn}(\Omega^{\check\rho})^{-\check\lambda}_D\simeq\Op^{RS}_G(D)_{\varpi(-\check\lambda-\check\rho)}.
\]
\end{proposition}

于是，若 $\check\lambda$ 满足命题 9.3.1 的条件，则 (9.3-6) 右端对应 $w=1$ 的层给出 Miura 变换 $\mu$ 在 $\Op^{RS}_G(D)_{\varpi(-\check\lambda-\check\rho)}$ 上的限制的一个截面。\transnote{卡点：(9.3-6) 右端按 $w\in W$ 分层的标号是 $\check\rho-w(\check\lambda+\check\rho)$。背景：Miura 变换给出的是 Cartan 侧联络 (9.3-4)，其留数为 $\check\lambda'-\check\rho$；(9.3-6) 固定 oper 侧留数 $\varpi(-\check\lambda-\check\rho)$，对应 Cartan 侧留数 $\check\rho-w(\check\lambda+\check\rho)$，$w=1$ 时恰为 $\check\rho-(\check\lambda+\check\rho)=-\check\lambda$。结论：$w=1$ 那一层就是 $\Conn(\Omega^{\check\rho})^{-\check\lambda}_D$，按命题 9.3.1 它与整个纤维同构，故给出 $\mu$ 在纤维上的一个截面。[来源：原书 p.266 (9.3-4)--(9.3-6) 与命题 9.3.1]}

% [Tutor Note: (9.3-6) 的右端是 $W$-分层，各层由 $\check\rho-w(\check\lambda+\check\rho)$ 标记；命题 9.3.1 说其中 $w=1$ 的那层（即 $\check\rho-(\check\lambda+\check\rho)=-\check\lambda$）与整个 $\Op^{RS}(D)_{\varpi(-\check\lambda-\check\rho)}$ 同构，因此 $\mu$ 在该留数纤维上有自然截面。这个“$w=1$ 层=全纤维”现象是 9.3 节全部计算的目标，也是第 10 章把 Miura 变换用于 Langlands 对应的技术基础。]

以下构造在后续有用。对 $D$ 上 $H$-丛 $\mathcal{F}_H$ 与整余权 $\check\lambda$，令 $\mathcal{F}_H(\check\lambda) = \mathcal{F}_H(\check\lambda\cdot 0)$，其中 $0$ 表示 $D$ 的原点（$D$ 的闭点），它是 $D$ 上这样的 $H$-丛：对任意特征 $\mu:H\to\C^\times$ 有
\[
\mathcal{F}_H(\check\lambda)\times_{\C^\times}\C_\mu = (\mathcal{F}_H\times_{\C^\times}\C_\mu)(\langle\mu,\check\lambda\rangle\cdot 0).
\]
注意丛 $\mathcal{F}_H(\check\lambda)$ 与 $\mathcal{F}_H$ 在 $D^\times$ 上的限制被典范等同。在此等同下，$\mathcal{F}_H(\check\lambda)$ 上的联络 $\nabla$ 成为 $\mathcal{F}_H$ 上的联络 $\check\lambda(t)^{-1}\nabla\check\lambda(t)$。

把 $D$ 上 $H$-丛 $\Omega^{\check\rho}(\check\lambda)$ 的联络空间记作 $\operatorname{Conn}(\Omega^{\check\rho}(\check\lambda))_D$。我们有自然嵌入
\[
\operatorname{Conn}(\Omega^{\check\rho}(\check\lambda))_D\subset\operatorname{Conn}(\Omega^{\check\rho})_{D^\times},
\]
它由把联络限制到 $D^\times$ 得到。$\operatorname{Conn}(\Omega^{\check\rho}(\check\lambda))_D$ 在 $\operatorname{Conn}(\Omega^{\check\rho})_{D^\times}$ 中的像与空间 $\operatorname{Conn}(\Omega^{\check\rho})^{\check\lambda}_D\subset\operatorname{Conn}(\Omega^{\check\rho})_{D^\times}$ 重合。确实，给定联络 $\nabla\in\operatorname{Conn}(\Omega^{\check\rho}(\check\lambda))_D$，联络 $\check\lambda(t)^{-1}\nabla\check\lambda(t)$ 属于 $\operatorname{Conn}(\Omega^{\check\rho})^{\check\lambda}_D$。

下文中，我们对这两个空间都使用记号 $\operatorname{Conn}(\Omega^{\check\rho})^{\check\lambda}_D$，同时记住上述两种不同的实现。

\subsection{整优余权的情形}

设 $\check\lambda$ 是 $G$ 的整优余权。圆盘 $D$ 上余权为 $\check\lambda$ 的\textbf{幂零 Miura $G$-oper}，按定义是四元组 $(\mathcal{F},\nabla,\mathcal{F}_B,\mathcal{F}'_B)$，其中 $(\mathcal{F},\nabla,\mathcal{F}_B)$ 是 $D$ 上余权为 $\check\lambda$ 的幂零 $G$-oper，而 $\mathcal{F}'_B$ 是 $\mathcal{F}$ 的另一个在 $\nabla_{t\partial_t} = t\nabla$ 下稳定的 $B$-约化。把余权 $\check\lambda$ 的幂零 Miura $G$-opers 构成的簇记作 $\mathrm{MOp}^{\check\lambda}_G$。

把余权为 $\check\lambda$ 的幂零 Miura $G$-oper 从 $D$ 限制到 $D^\times$，得到 8.2.1 节所定义的 $D^\times$ 上的 Miura oper。把 $D^\times$ 上 Miura opers 的集合记作 $\mathrm{MOp}_G(D^\times)$。我们有自然的遗忘映射 $\mathrm{MOp}^{\check\lambda}_G\to\Op^{\mathrm{nilp},\check\lambda}_G$，它使下图成为 Cartesian 图表：
\[
\begin{CD}
\mathrm{MOp}^{\check\lambda}_G @>>> \mathrm{MOp}_G(D^\times)\\
@VVV @VVV\\
\Op^{\mathrm{nilp},\check\lambda}_G @>>> \Op_G(D^\times)
\end{CD}
\]

下一个引理的证明属于 V. Drinfeld（见 [2006c]，引理 3.2.1）。回忆 8.2.1 节：Miura oper $(\mathcal{F},\nabla,\mathcal{F}_B,\mathcal{F}'_B)$ 称为\textbf{泛的}，如果 $B$-约化 $\mathcal{F}_B$ 与 $\mathcal{F}'_B$ 处于泛相对位置。

\begin{lemma}
去心圆盘 $D^\times$ 上任一 Miura oper 都是泛的。
\end{lemma}

\begin{proof}
在圆盘 $D$ 上取坐标 $t$，并平凡化 $B$-丛 $\mathcal{F}_B$。于是把 $\mathcal{F}$ 在 $D^\times$ 上的截面与 $G((t))$ 等同，把约化 $\mathcal{F}'_B$ 与右陪集 $g\cdot B((t))\subset G((t))$ 等同，把联络 $\nabla$ 与算子
\begin{equation}
\partial_t + A(t),\qquad A(t) = \sum_{i=1}^{\ell}\psi_i(t)f_i + \mathbf{v}(t)
\tag{9.3-7}
\end{equation}
等同，其中 $\psi_i(t)\in\C((t))^\times$，$\mathbf{v}(t)\in\fb[[t]]$。

有 Bruhat 分解 $G((t)) = \bigsqcup_{w\in W}B((t))wB((t))$。因此可写 $g = g_1wg_2$，其中 $w\in W$，$g_1,g_2\in B((t))$。于是 Borel 子代数 $\fb_{\mathcal{F}'_B}\subset\fg_{\mathcal{F}'_B} = \fg_{\mathcal{F}_B}$ 在 $D^\times$ 上的截面空间等于
\[
\fb' = g\fb((t))g^{-1} = g_1(w\fb((t))w^{-1})g_1^{-1}.
\]
引理的断言等价于：此式中的元素 $w$ 是外尔群 $W$ 的最长元素 $w_0$。

为看出这一点，考虑 Cartan 子代数 $\fh' = g_1\fh((t))g_1^{-1}$。则由于 $g_1\in B((t))$，有 $\fh'\subset\fb'\cap\fb((t))$。关于 $\fh'$ 有如下 $\fg((t))$ 的 Cartan 分解：
\[
\fg((t)) = \fh'\oplus\bigoplus_{\alpha\in\Delta}\fg'_\alpha,\qquad \fg'_\alpha = g_1\fg_\alpha g_1^{-1},
\]
其中 $\fg_\alpha$ 是 $\fg$ 中对应 $\alpha\in\Delta$ 的根子空间。把 (9.3-7) 中的 $A(t)$ 关于此分解展开：
\[
A(t) = A_0 + \sum_{\alpha\in\Delta}A_\alpha.
\]
则由 oper 条件有 $A_{-\alpha_i}\ne 0$。另一方面，由于 $\mathcal{F}'_B$ 被 $\nabla$ 保持，有 $A(t)\in\fb'$。这意味着相对于 $\fb'$，根 $-\alpha_i$ 都是正根。因此相对于 $\fb((t))$ 的一切负根都是相对于 $\fb'$ 的正根。这当且仅当 $\fb' = g_1\fb_-g_1^{-1}$，其中 $\fb_- = w_0\fb w_0^{-1}$，$g_1\in B((t))$。\transnote{卡点：由 $A_{-\alpha_i}\ne0$ 何以断言 $w=w_0$。背景：$\fb'$ 被 $\nabla$ 保持，故 $A(t)\in\fb'$；而 oper 条件要求每个简单根方向的 $A_{-\alpha_i}$ 非零。于是相对 $\fb'$，全部 $-\alpha_i$ 都是正根，$\fb'$ 的根系必含所有 $-\alpha_i$，即它与 $\fb$ 在 Bruhat 分解中处于最远的位置。结论：$\fb'=g_1\fb_-g_1^{-1}$（$\fb_-=w_0\fb w_0^{-1}$），对应 $w=w_0$。[来源：原书 p.268]}
\end{proof}

% [Tutor Note: 证明的关键判据是 $A_{−\alpha_i}\ne 0$（注意下标是 $-\alpha_i$）。oper 条件保证每个单根方向的系数非零；若 $\mathcal{F}'_B$ 与 $\mathcal{F}_B$ 处于泛位置，则所有负根相对于 $\fb'$ 都变正，这迫使 $w$ 为最长元素 $w_0$。抄写时容易把 $A_{-\alpha_i}$ 写成 $A_{\alpha_i}$，那样论证方向恰好相反。]

与引理 8.2.2 相同可证：$D^\times$ 上泛 Miura opers 的集合与 $H$-丛 $\Omega^{\check\rho}$ 在 $D^\times$ 上的联络空间 $\operatorname{Conn}(\Omega^{\check\rho})_{D^\times}$ 之间存在双射。再由引理 9.3.2 得到：

\begin{lemma}
存在双射
\[
\mathrm{MOp}_G(D^\times)\simeq\operatorname{Conn}(\Omega^{\check\rho})_{D^\times}.
\]
\end{lemma}

于是，任何幂零 Miura oper 限制到去心圆盘后都变成泛的，从而定义了 $\Omega^{\check\rho}$ 上的一个联络。现在我们的目标是理解这个联络的留数与原点 $0\in D$ 处 $B$-约化 $\mathcal{F}_B$ 与 $\mathcal{F}'_B$ 的相对位置之间的关系。

设 $\mathcal{N}\subset\fg$ 是 $\fg$ 中幂零元素构成的锥，$\widetilde{\mathcal{N}}$ 是 Springer 簇，即对 $(x,\fb')$（$x\in\mathcal{N}$，$\fb'$ 是 $\fg$ 中含 $x$ 的 Borel 子代数）构成的簇。设 $\tilde{\mathfrak{n}}$ 是 $\widetilde{\mathcal{N}}$ 到 $\mathfrak{n}\subset\mathcal{N}$ 的限制。于是
\[
\tilde{\mathfrak{n}} = \{(x,\fb')\mid x\in\mathfrak{n},\ x\in\fb'\}.
\]
回忆 9.2.3 节定义了余权为 $\check\lambda$ 的幂零 oper $\chi$ 的留数 $\operatorname{Res}_{\check\lambda}(\chi)$。

\begin{lemma}
余权为 $\check\lambda$ 的幂零 oper $\chi$ 的约化 $\mathcal{F}'_B$ 由它在原点的纤维唯一决定。它必须在底层幂零 oper 的留数 $\operatorname{Res}_{\check\lambda}(\chi)$ 下稳定。
\end{lemma}

为看出这一点，回忆 $\mathcal{F}'_B$ 可以通过在向量丛 $V^\nu_{\mathcal{F}} = \mathcal{F}\times_G V^\nu$ 中选取线子丛 $\mathcal{L}^\nu$（其中 $V^\nu$ 是 $G$ 的最高权为 $\nu$ 的不可约表示）并满足 Plücker 关系来完全描述（例如见 [Frenkel et al. 2001] 第 2.1.2 节）。% [Tutor Note: 这里“$\mathcal{F}'_B$ 由线子丛 $\mathcal{L}^\nu$ 满足 Plücker 关系完全描述”用的是把 $B$-约化翻译成各不可约表示中线丛的标准办法：$\mathcal{F}'_B$ 对应每个 $V^\nu$ 中的最高权线子丛。引理 9.3.5 的“稳定”正好等于 Plücker 关系在 $t=0$ 处有解，因而 引理 9.3.4 是它的推论。]
于是上述断言是下述引理的推论：

\begin{lemma}
设给定有限维向量空间 $V$ 上带正则奇点的一阶微分方程
\[
(\partial_t + A(t))\Psi(t) = 0,\qquad A(t) = \sum_{i\ge -1}A_it^i.
\]
设 $A_{-1}$ 是 $V$ 上的幂零算子。则该方程有满足初值条件 $\Psi(0)=v\in V$ 的解 $\Psi(t)\in V[[t]]$，当且仅当 $A_{-1}(v)=0$。
\end{lemma}

\begin{proof}
把 $\Psi(t) = \sum_{j\ge 0}\Psi_jt^j$ 展开。则上述方程等价于线性方程组
\begin{equation}
n\Psi_n + \sum_{i+j=n-1}A_i(\Psi_j) = 0.
\tag{9.3-8}
\end{equation}
其中第一个方程是 $A_{-1}(\Psi_0)=0$，所以这确实是必要条件。若此条件满足，则 $n>0$ 时的 $\Psi_n$ 可递推求出，因为 $A_{-1}$ 的幂零性蕴含算子 $(n\operatorname{Id}+A_{-1})$（$n>0$）可逆。
\end{proof}

这意味着有自然同构
\begin{equation}
\mathrm{MOp}^{\check\lambda}_G\simeq\Op^{\mathrm{nilp},\check\lambda}_G\underset{\mathfrak{n}/B}{\times}\tilde{\mathfrak{n}}/B,
\tag{9.3-9}
\end{equation}
其中映射 $\Op^{\mathrm{nilp},\check\lambda}_G\to\mathfrak{n}/B$ 是留数映射 $\operatorname{Res}_{\check\lambda}$。于是 $\mathrm{MOp}^{\check\lambda}_G$ 在 $\chi\in\Op^{\mathrm{nilp},\check\lambda}_G$ 上的纤维是 $\operatorname{Res}_{\check\lambda}(\chi)$ 的（既约）Springer 纤维：
\begin{equation}
\mathrm{Sp}_{\operatorname{Res}_{\check\lambda}(\chi)} = \{\fb'\in G/B\mid \operatorname{Res}_{\check\lambda}(\chi)\in\fb'\}.
\tag{9.3-10}
\end{equation}

给定 $w\in W$，我们说 Borel 子群 $B'$（或相应的李代数 $\fb'$）与 $B$（相应地 $\fb$）处于\textbf{相对位置 $w$}，如果 $B'$（相应地 $\fb'$）所对应的旗簇 $G/B$ 中的点落在 $B$-轨道 $Bw^{-1}w_0B$ 中。把 $\widetilde{\mathcal{N}}$ 中由这样的对 $(x,\fb')$ 构成的子簇记作 $\widetilde{\mathcal{N}}_w$：其中 $\fb'$ 与我们固定的 Borel 子代数 $\fb$ 处于相对位置 $w$。相应地，把 $\widetilde{\mathcal{N}}_w$ 到 $\mathfrak{n}$ 的限制记作 $\tilde{\mathfrak{n}}_w$。群 $B$ 自然作用在 $\tilde{\mathfrak{n}}$ 上并保持 $\tilde{\mathfrak{n}}_w$。

把 $\mathrm{MOp}^{\check\lambda}_G$ 中那些 $\mathcal{F}'_B$ 在原点的纤维与 $B$ 处于相对位置 $w$ 的 Miura opers 构成的子簇记作 $\mathrm{MOp}^{\check\lambda,w}_G$。于是有等同
\begin{equation}
\mathrm{MOp}^{\check\lambda,w}_G = \Op^{\mathrm{nilp},\check\lambda}_G\underset{\mathfrak{n}/B}{\times}\tilde{\mathfrak{n}}_w/B.
\tag{9.3-11}
\end{equation}

\subsection{联络与 Miura opers}

按命题 9.2.1，对任意整优余权 $\check\lambda$ 有同构
\[
\Op^{RS}_G(D)_{-\varpi(\check\lambda+\check\rho)}\simeq\Op^{\mathrm{nilp},\check\lambda}_G.
\]
因此，态射 $\mu^{\mathrm{RS}}$（见公式 (9.3-4)）在子空间
\[
\operatorname{Conn}(\Omega^{\check\rho})^{-w(\check\lambda+\check\rho)+\check\rho}_D\subset\operatorname{Conn}(\Omega^{\check\rho})^{\mathrm{RS}}_D
\]
上的限制给出态射
\[
\mu_{-w(\check\lambda+\check\rho)+\check\rho}: \operatorname{Conn}(\Omega^{\check\rho})^{-w(\check\lambda+\check\rho)+\check\rho}_D\to\Op^{\mathrm{nilp},\check\lambda}_G.
\]
现在，由 Miura 变换 $\mu$ 的显式构造（见 9.3.1 节）可知：对任何余权为 $\check\lambda$ 的幂零 oper $\chi$（它属于映射 $\mu_{-w(\check\lambda+\check\rho)+\check\rho}$ 的像），Borel 约化
\[
\mathcal{F}'_B = \Omega^{\check\rho}\underset{H}{\times}w_0B
\]
在 $t\nabla$ 下稳定。因此映射 $\mu_{-w(\check\lambda+\check\rho)+\check\rho}$ 可以典范地提升为映射
\begin{equation}
\tilde\mu_{-w(\check\lambda+\check\rho)+\check\rho}: \operatorname{Conn}(\Omega^{\check\rho})^{-w(\check\lambda+\check\rho)+\check\rho}_D\to\mathrm{MOp}^{\check\lambda}_G.
\tag{9.3-12}
\end{equation}
这个映射使下图成为 Cartesian 图表：
\[
\begin{CD}
\operatorname{Conn}(\Omega^{\check\rho})^{-w(\check\lambda+\check\rho)+\check\rho}_D @>>> \operatorname{Conn}(\Omega^{\check\rho})_{D^\times}\\
@V{\tilde\mu_{-w(\check\lambda+\check\rho)+\check\rho}}VV @VVV\\
\mathrm{MOp}^{\check\lambda}_G @>>> \mathrm{MOp}_G(D^\times)
\end{CD}
\]
其中右竖直映射是引理 9.3.3 中建立的双射。下述结果在 [Frenkel and Gaitsgory 2006c]，Cor. 3.6.5 中证明（对 $\check\lambda=0$）。

\begin{theorem}
映射 $\tilde\mu_{-w(\check\lambda+\check\rho)+\check\rho}$ 是同构
\[
\operatorname{Conn}(\Omega^{\check\rho})^{-w(\check\lambda+\check\rho)+\check\rho}_D\simeq\mathrm{MOp}^{\check\lambda,w}_G\subset\mathrm{MOp}^{\check\lambda}_G.
\]
\end{theorem}

\begin{proof}
设给定 $\mathrm{MOp}^{\check\lambda}_G$ 中一点，即一个余权为 $\check\lambda$、配备在 $t\nabla$ 下稳定的 Borel 约化 $\mathcal{F}'_B$ 的幂零 oper。由于我们的幂零 Miura oper 在 $D^\times$ 上是泛的（引理 9.3.2），由引理 8.2.1 可知 $H$-丛 $\mathcal{F}'_B/N$ 在 $D^\times$ 上同构于 $w_0^*(\mathcal{F}_B/N) = w_0^*(\Omega^{\check\rho})$。于是得到 $w_0^*(\Omega^{\check\rho})$ 上的联络，从而得到 $\Omega^{\check\rho}$ 在 $D^\times$ 上的联络。把后者记作 $\overline{\nabla}$。显然 $\overline{\nabla}$ 有正则奇点，故
\[
\overline{\nabla} = \partial_t - \frac{1}{t}(\check\mu+t(\dots)),\qquad \check\mu\in\fh.
\]
显然，由 $\overline{\nabla}$ 经引理 9.3.3 的双射得到的 $D^\times$ 上 Miura oper $(\mathcal{F},\nabla,\mathcal{F}_B,\mathcal{F}'_B)$ 正是我们出发时的 Miura oper。因此相应的 oper 具有形式
\[
\partial_t + \frac{1}{t}(p_{-1}-\check\rho-\check\mu+t(\dots)).
\]
于是它在 $\fh/W$ 中的留数等于 $\varpi(-\check\mu-\check\rho)$。但按假设这个 oper 的留数等于 $\varpi(-\check\lambda-\check\rho)$。因此存在 $y\in W$ 使 $-\check\rho-\check\mu = y(-\check\lambda-\check\rho)$，即 $\check\mu = y(\check\lambda+\check\rho)-\check\rho$。\transnote{卡点：$y\in W$ 从何而来。背景：证明第一步由引理 8.2.1（$B$-约化在 $D^\times$ 上处于泛位置）得 $\mathcal{F}'_B/N\simeq w_0^*(\mathcal{F}_B/N)=w_0^*(\Omega^{\check\rho})$，把点化为 $\Omega^{\check\rho}$ 上的联络 $\overline\nabla$。留数只在 $\fh/W$ 中有定义，故两侧留数类相等只给出 $-\check\rho-\check\mu=y(-\check\lambda-\check\rho)$，$y$ 即此处的 $W$-平移。结论：$y$ 正是 (9.3-6) 右端分层的参数，也是定理 9.3.6 的像只落在 $\mathrm{MOP}^{\check\lambda,w}_G$ 而不落在整个 $\mathrm{MOP}^{\check\lambda}_G$ 的原因。[来源：原书 p.270--271、引理 8.2.1]}

于是得到：引理 9.3.3 的双射限制到 $\mathrm{MOp}^{\check\lambda}_G$ 上给出双射
\[
\mathrm{MOp}^{\check\lambda}_G\simeq\bigsqcup_{y\in W}\operatorname{Conn}(\Omega^{\check\rho})^{-y(\check\lambda+\check\rho)+\check\rho}_D.
\]
我们需要证明：在此双射下，$\mathrm{MOp}^{w,\check\lambda}_G$ 恰好对应
\[
\operatorname{Conn}(\Omega^{\check\rho})^{-w(\check\lambda+\check\rho)+\check\rho}_D.
\]
回忆在 $D^\times$ 上有 $H$-丛的同构 $\mathcal{F}_H|_{D^\times}\simeq w_0^*(\mathcal{F}'_H)|_{D^\times}$，其中 $\mathcal{F}_H = \mathcal{F}_B/N$，$\mathcal{F}'_H = \mathcal{F}'_B/N$。注意：若有两个 $D$ 上的 $H$-丛 $\mathcal{F}_H$ 与 $\mathcal{F}''_H$，则任一在 $D^\times$ 上的同构 $\mathcal{F}_H|_{D^\times}\simeq\mathcal{F}''_H|_{D^\times}$ 都扩张为 $D$ 上的同构 $\mathcal{F}_H\simeq\mathcal{F}''_H(\check\mu)$，其中 $\check\mu$ 是某个整余权。按引理 4.2.1 证明中的方式考虑幂零 oper 的变换性质，可得
\[
\mathcal{F}_H = \mathcal{F}_B/N\simeq\Omega^{\check\rho}(-\check\lambda).
\]

若我们能证明同构 $\mathcal{F}_H|_{D^\times}\simeq w_0^*(\mathcal{F}'_H)|_{D^\times}$ 扩张为 $D$ 上的同构
\begin{equation}
\mathcal{F}'_H\simeq w_0^*\bigl(\mathcal F_H((\check\lambda+\check\rho)-w(\check\lambda+\check\rho))\bigr),
\tag{9.3-13}
\end{equation}
则定理的断言随之成立，因为此时
\[
\mathcal{F}'_H\simeq w_0^*(\Omega^{\check\rho}(\check\rho-w(\check\lambda+\check\rho))),
\]
于是 $D^\times$ 上 $\Omega^{\check\rho}$ 上的诱导联络落在
\[
\operatorname{Conn}(\Omega^{\check\rho}(\check\rho-w(\check\lambda+\check\rho)))_D\simeq\operatorname{Conn}(\Omega^{\check\rho})^{\check\rho-w(\check\lambda+\check\rho)}_D
\]
的像中，这正是所要的。

让我们这样平凡化 $\mathcal{F}$ 在原点的纤维：使 $\mathcal{F}_B$ 与 $B\subset G$ 等同。不妨设在平凡化下 $\mathcal{F}'_{B,0}$ 变成 $w^{-1}w_0B$。需要证明下述事实。对整优权 $\nu$，设 $V^\nu$ 是 $\fg$ 的最高权为 $\nu$ 的有限维表示。对 $w\in W$，在 $V^\nu$ 的权空间 $w(\nu)$ 的一维子空间中取非零向量 $v_w$。按假设，$\mathcal{F}'_B$ 被 $t\nabla$ 保持。因此由引理 9.3.4，$\nabla$ 的留数稳定向量 $v_{w^{-1}w_0}$。考虑相伴向量丛 $V^\nu_{\mathcal{F}} = \mathcal{F}\times_G V^\nu$ 以及由 $\nabla$ 诱导的联络。由引理 9.3.5 可知，这个平坦丛存在唯一的单值水平截面，它在 $t=0$ 处的值等于 $v_{w^{-1}w_0}$。把 $\Psi(t)$ 按不同权分量展开。设 $\phi_w(t)v_{w_0}$（$\phi_w(t)\in\C[[t]]$）是 $\Psi(t)$ 中最低权 $w_0(\nu)$ 的分量。由引理 9.3.2，$\mathcal{F}'_B$ 与 $\mathcal{F}_B$ 在 $D^\times$ 上处于泛位置。这意味着 $\phi_w(t)\ne 0$。设 $n_w$ 是 $\phi_w(t)$ 在 $t=0$ 处的消没阶数。

我们需要证明：对每个整优权 $\nu$ 都有
\begin{equation}
n_w = \langle w_0(\nu), w(\check\lambda+\check\rho)-(\check\lambda+\check\rho)\rangle,
\tag{9.3-14}
\end{equation}
因为这等价于公式 (9.3-13)。

为计算 $\phi_w(t)$ 的消没阶数，回忆余权为 $\check\lambda$ 的幂零 oper 底下的联络具有形式
\begin{equation}
\nabla = \partial_t + \sum_{i=1}^{\ell}t^{\langle\alpha_i,\check\lambda\rangle}f_i + \mathbf{v}(t) + \frac{v}{t},
\tag{9.3-15}
\end{equation}
其中 $\mathbf{v}(t)\in\fb[[t]]$，$v\in\fn$。$V^\nu_{\mathcal{F}}$ 的所求水平截面可通过在 $V^\nu[[t]]$ 中解方程 $\nabla\Psi(t)=0$ 得到。像引理 9.3.5 的证明那样递推求解这个方程。于是需要解方程 (9.3-8)，其中右端是 $\fb$ 中算子的线性组合（这些算子保持或提高主分次），以及算子 $f_i$（$i=1,\dots,\ell$，它们把主分次降低 1）。因此显然，$\Psi(t)$ 中正比于 $v_{w_0}$ 的分量 $\phi_w(t)v_{w_0}$ 的首项，是通过把算子 $f_i$（$i=1,\dots,\ell$）作用于 $v_{w^{-1}w_0}$ 得到的。故它与联络
\[
\nabla = \partial_t + \sum_{i=1}^{\ell}t^{\langle\alpha_i,\check\lambda\rangle}\psi_i(t)f_i
\]
的情形相同。当对这个联络求解递推关系 (9.3-8) 时，每作用一次 $f_i$ 就把 $\Psi(t)$ 的 $t$-次数增加 $\langle\alpha_i,\check\lambda\rangle+1 = \langle\alpha_i,\check\lambda+\check\rho\rangle$。这是 $\alpha_i$ 的“$(\check\lambda+\check\rho)$-次数”。因此，要从 $v_{w^{-1}w_0}$（$\Psi(t)$ 的初值条件）“到达” $v_{w_0}$，需要把 $t$-次数增加向量 $v_{w^{-1}w_0}$ 与 $v_{w_0}$ 的 $(\check\lambda+\check\rho)$-次数之差，即增加
\[
\langle w^{-1}w_0(\nu),\check\lambda+\check\rho\rangle-\langle w_0(\nu),\check\lambda+\check\rho\rangle = \langle w_0(\nu),w(\check\lambda+\check\rho)-(\check\lambda+\check\rho)\rangle,
\]
这正是公式 (9.3-14) 给出的 $n_w$。于是得到 $\phi_w(t) = c_wt^{n_w}$，其中 $c_w\in\C$。若 $c_w=0$，则解 $\Psi(t)$ 的整个 $v_{w_0}$-分量都等于 0。但这将意味着相应的 Borel 约化 $\mathcal{F}'_B$ 与 $\mathcal{F}_B$ 在 $D^\times$ 上不处于泛位置，与引理 9.3.2 矛盾。因此 $c_w\ne 0$，从而对任意 oper (9.3-15)，$\phi_w(t)$ 的首项 $t$-次数确实等于 $n_w$。证明完毕。\transnote{卡点：定理 9.3.6 的证明在此收尾，其中三处跳步。背景：(i) 证明只追踪 $v_{w_0}$-分量的首项次数，原书未展开其余项为何不影响首项；(ii) 每次作用 $f_i$ 恰好把 $t$-次数提高 $\langle\alpha_i,\check\lambda\rangle+1=\langle\alpha_i,\check\lambda+\check\rho\rangle$；(iii) 只需对一切整优权 $\nu$ 验证 (9.3-14)，因为 $V^\nu$ 的权集合决定 $W$-平移，(9.3-14) 与 (9.3-13) 等价。结论：把上述次数累加即得 (9.3-14)，从而 $\mathcal{F}'_H\simeq w_0^*(\mathcal{F}_H(\dots))$ 延拓到 $D$ 上。[来源：原书 p.272--273；Frenkel and Gaitsgory 2006c, Cor. 3.6.5（原书 p.271 自引）]}
\end{proof}

% [Tutor Note: 这段证明的技术核心是“$(\check\lambda+\check\rho)$-次数”：$f_i$ 每作用一次把 $t$-次数提高 $\langle\alpha_i,\check\lambda+\check\rho\rangle$，于是从最低权向量 $v_{w^{-1}w_0}$ 到最高权向量 $v_{w_0}$ 所需的次数恰由 (9.3-14) 给出。这里的 $w^{-1}w_0(\nu)$ 应读作 $(w^{-1}w_0)(\nu)$，阅读时不要误当成 $w^{-1}$ 作用后再乘 $w_0\nu$。]

结合同构 (9.3-12) 与 (9.3-11)，得到同构
\begin{equation}
\operatorname{Conn}(\Omega^{\check\rho})^{\check\rho-w(\check\lambda+\check\rho)}_D\simeq\Op^{\mathrm{nilp},\check\lambda}_G\underset{\mathfrak{n}/B}{\times}\tilde{\mathfrak{n}}_w/B.
\tag{9.3-16}
\end{equation}
于是得到 $\C$-点集合之间的双射
\begin{equation}
\bigsqcup_{w\in W}\operatorname{Conn}(\Omega^{\check\rho})^{\check\rho-w(\check\lambda+\check\rho)}_D\simeq\Op^{\mathrm{nilp},\check\lambda}_G\underset{\mathfrak{n}/B}{\times}\tilde{\mathfrak{n}}/B.
\tag{9.3-17}
\end{equation}
虽然 (9.3-16) 是簇的同构，但 (9.3-17) 不是，因为左端是若干层的无交并，而它们在右端被“粘合”在一起。

回忆按公式 (9.3-6)，(9.3-17) 的左端正是 $\Op^{\check\lambda}_G\subset\Op_G(D^\times)$ 在 Miura 变换 $\mu$ 下的原像。于是我们看到，在同构 (9.3-17) 下，自然投影
\[
\Op^{\mathrm{nilp},\check\lambda}_G\underset{\mathfrak{n}/B}{\times}\tilde{\mathfrak{n}}/B\to\Op^{\mathrm{nilp},\check\lambda}_G
\]
到第一个因子的投影，与 Miura 变换 $\mu$（见 (9.3-3)）在 $\Op^{\mathrm{nilp},\check\lambda}_G$ 上的限制一致。这意味着，纤维 $\mu^{-1}(\chi)$ 的点集与 Springer 纤维 $\operatorname{Sp}_{\operatorname{Res}_{\check\lambda}(\chi)}$ 的点集一一对应（见 (9.3-10)）。
% （本单元止于 PDF p.283 页末；下接 ch09-2.tex，该页 p.284 上半页的 §9.3 续文从“与同构 (9.3-17) 的情形一样”起。）
